\documentclass[11pt,a4paper]{article}

\usepackage{fontspec}
\usepackage[margin=2.4cm]{geometry}
\usepackage{xcolor}
\usepackage{booktabs}
\usepackage{amsmath}
\usepackage{enumitem}
\usepackage{listings}
\usepackage{tcolorbox}
\usepackage{titlesec}
\usepackage{fancyhdr}
\usepackage[hidelinks]{hyperref}

\setmainfont{texgyrepagella}[Extension=.otf, UprightFont=*-regular,
  BoldFont=*-bold, ItalicFont=*-italic, BoldItalicFont=*-bolditalic]
\setsansfont{texgyreheros}[Extension=.otf, UprightFont=*-regular,
  BoldFont=*-bold, ItalicFont=*-italic, BoldItalicFont=*-bolditalic]
\setmonofont{lmmono10}[Extension=.otf, UprightFont=*-regular,
  Scale=0.9]

\definecolor{accent}{HTML}{2A7D78}
\definecolor{ink}{HTML}{1E2633}
\definecolor{muted}{HTML}{64708A}
\definecolor{codebg}{HTML}{F1F3F6}
\definecolor{boxbg}{HTML}{EAF4F3}

\color{ink}
\setlength{\parindent}{0pt}
\setlength{\parskip}{0.55em}
\linespread{1.08}
\emergencystretch=3em

\titleformat{\section}{\sffamily\Large\bfseries\color{accent}}{\thesection}{0.6em}{}
\titleformat{\subsection}{\sffamily\large\bfseries}{\thesubsection}{0.6em}{}
\titlespacing*{\section}{0pt}{1.6em}{0.6em}
\titlespacing*{\subsection}{0pt}{1.2em}{0.4em}

\pagestyle{fancy}
\fancyhf{}
\renewcommand{\headrulewidth}{0pt}
\fancyfoot[L]{\small\color{muted}Nordhjem store \textperiodcentered{} Step 1}
\fancyfoot[R]{\small\color{muted}\thepage}

\lstset{
  basicstyle=\ttfamily\small,
  backgroundcolor=\color{codebg},
  frame=none,
  xleftmargin=0.8em, xrightmargin=0.8em,
  aboveskip=0.8em, belowskip=0.8em,
  columns=fullflexible, keepspaces=true,
  showstringspaces=false,
  keywordstyle=\color{accent}\bfseries,
  commentstyle=\color{muted}\itshape,
  literate={×}{{$\times$}}1 {·}{{$\cdot$}}1
           {æ}{{\ae}}1 {ø}{{\o}}1 {å}{{\aa}}1 {Æ}{{\AE}}1
}

\newtcolorbox{task}[1]{colback=boxbg, colframe=accent, boxrule=0.6pt,
  arc=3pt, left=8pt, right=8pt, top=6pt, bottom=6pt,
  fonttitle=\sffamily\bfseries, title=#1}
\newtcolorbox{note}{colback=codebg, colframe=codebg, arc=3pt,
  left=8pt, right=8pt, top=5pt, bottom=5pt}

\newcommand{\code}[1]{\texttt{#1}}

\begin{document}

{\sffamily
{\small\color{muted} NORDHJEM STORE \textperiodcentered{} AGENTS AND RAG ON GOOGLE CLOUD}\\[0.4em]
{\LARGE\bfseries Step 1: semantic search, by hand}\\[0.6em]
}

\textbf{Goal.} A customer types \emph{``a couch my guests can sleep on''}
in the Nordhjem search box and gets sofas that fold out into a bed, even
though no product text contains the word ``couch''. You write this
search yourself, in \code{api/search.py}, with the embedding model and
NumPy only, so every step is visible. When it passes the check, the live
page switches from ``keyword placeholder'' to ``your search''.

What is already done for you (plumbing): the 500 products are in
Firestore, each with its vector of 768 numbers. Your job is everything
that happens when a customer asks.

\tableofcontents

%--------------------------------------------------------------------
\section{The problem with keyword search}

The page runs a keyword search today (\code{api/keyword\_search.py}). It
counts how many words of the question appear in a product's text. These
are real results from it:

\begin{itemize}[leftmargin=1.4em]
  \item \emph{``a couch my guests can sleep on''}: 2 sofas and 3 rugs,
        all with the same score $0.43$. Why $0.43$: the question has 7
        words, 3 of them are found, and $3 / 7 \approx 0.43$. The 3 words
        found are small common ones, not ``couch''.
  \item \emph{``couch''}: nothing at all. No product text contains that
        word; the catalogue says ``sofa''.
\end{itemize}

The catalogue does have the right products, for example \emph{``2-seat
sofa in grey fabric. Folds out into a bed for guests''}. Keyword search
matches \textbf{letters}. We want to match \textbf{meaning}. The
reference solution of this step returns 5 such sofa beds for the same
question, with scores around $0.71$.

%--------------------------------------------------------------------
\section{What an embedding is}

An \textbf{embedding model} reads a text and returns a list of numbers,
called a \textbf{vector}. Here the model is \code{gemini-embedding-001}
on Vertex AI, and we ask for 768 numbers per text. The model was trained
so that texts with similar meaning get vectors that point in a similar
direction, in any language.

A worked example, with made-up vectors of only 3 numbers so you can
follow by hand:

\begin{center}
\begin{tabular}{@{}llc@{}}
\toprule
 & Text & Vector \\
\midrule
A & ``sofa that folds out into a bed'' & $[0.9,\ 0.1,\ 0.4]$ \\
B & ``a couch my guests can sleep on'' & $[0.8,\ 0.2,\ 0.5]$ \\
C & ``wool rug, 160x230 cm''           & $[0.1,\ 0.9,\ 0.1]$ \\
\bottomrule
\end{tabular}
\end{center}

A and B point almost the same way; C points elsewhere. So the search
question becomes: \emph{which product vectors point most like the
question's vector?}

%--------------------------------------------------------------------
\section{Measuring ``points the same way'': cosine similarity}

We use the three vectors from Section~2:
\[
A = [0.9,\ 0.1,\ 0.4] \qquad B = [0.8,\ 0.2,\ 0.5] \qquad C = [0.1,\ 0.9,\ 0.1]
\]
The \textbf{cosine similarity} of two vectors $a$ and $b$ is
\[
\cos(a, b) \;=\; \frac{a \cdot b}{\lVert a \rVert \times \lVert b \rVert}
\]
It has two ingredients: the dot product $a \cdot b$ on top, and the
lengths $\lVert a \rVert$ and $\lVert b \rVert$ below. We compute each
one, step by step. Results are rounded to 3 decimals.

\subsection{Ingredient 1: the dot product}

\textbf{Definition.} Multiply the first numbers of both vectors, then the
second numbers, then the third, and add the products:
\[
a \cdot b = a_1 b_1 + a_2 b_2 + a_3 b_3
\]

\textbf{A and B.}
\[
A \cdot B = 0.9 \times 0.8 \;+\; 0.1 \times 0.2 \;+\; 0.4 \times 0.5
          = 0.72 + 0.02 + 0.20 = 0.94
\]

\textbf{A and C.}
\[
A \cdot C = 0.9 \times 0.1 \;+\; 0.1 \times 0.9 \;+\; 0.4 \times 0.1
          = 0.09 + 0.09 + 0.04 = 0.22
\]

\textbf{B and C.}
\[
B \cdot C = 0.8 \times 0.1 \;+\; 0.2 \times 0.9 \;+\; 0.5 \times 0.1
          = 0.08 + 0.18 + 0.05 = 0.31
\]

\subsection{Ingredient 2: the length (L2 norm)}

\textbf{Definition.} Square each number of the vector, add the squares,
and take the square root:
\[
\lVert a \rVert = \sqrt{a_1^2 + a_2^2 + a_3^2}
\]
This is the \textbf{Euclidean norm}, also called the \textbf{L2 norm}
(the ``2'' is the power in the squares). It is the straight-line length
of the arrow. In 2 dimensions it is Pythagoras: the vector $[3, 4]$ has
length $\sqrt{3^2 + 4^2} = \sqrt{9 + 16} = \sqrt{25} = 5$. In NumPy,
\code{np.linalg.norm(v)} computes it.

\textbf{Length of A.}
\[
\lVert A \rVert = \sqrt{0.9^2 + 0.1^2 + 0.4^2}
  = \sqrt{0.81 + 0.01 + 0.16} = \sqrt{0.98} \approx 0.990
\]

\textbf{Length of B.}
\[
\lVert B \rVert = \sqrt{0.8^2 + 0.2^2 + 0.5^2}
  = \sqrt{0.64 + 0.04 + 0.25} = \sqrt{0.93} \approx 0.964
\]

\textbf{Length of C.}
\[
\lVert C \rVert = \sqrt{0.1^2 + 0.9^2 + 0.1^2}
  = \sqrt{0.01 + 0.81 + 0.01} = \sqrt{0.83} \approx 0.911
\]

\subsection{Putting it together: the cosine}

\textbf{A and B} (the sofa bed and the couch for guests):
\[
\cos(A, B) = \frac{A \cdot B}{\lVert A \rVert \times \lVert B \rVert}
  = \frac{0.94}{0.990 \times 0.964}
  = \frac{0.94}{0.955}
  \approx \mathbf{0.985}
\]

\textbf{A and C} (the sofa bed and the rug):
\[
\cos(A, C) = \frac{0.22}{0.990 \times 0.911}
  = \frac{0.22}{0.902}
  \approx \mathbf{0.244}
\]

\textbf{B and C} (the couch for guests and the rug):
\[
\cos(B, C) = \frac{0.31}{0.964 \times 0.911}
  = \frac{0.31}{0.879}
  \approx \mathbf{0.353}
\]

\textbf{Reading the result.} The question is B. Its scores are $0.985$
with the sofa bed A and $0.353$ with the rug C. Ranked from high to low:
A first, C second. That ranking is the search result.

The cosine is always between $-1$ and $1$: $1$ means exactly the same
direction, $0$ unrelated (a right angle), $-1$ opposite.

\subsection{The shortcut: normalise once, then only dot products}

\textbf{Normalising} means dividing every number of a vector by the
vector's length. The result points the same way and has length exactly
1. We write it with a hat: $\hat{A}$.

\textbf{Normalising A.} The length of $A$ is $0.990$ (more exactly
$0.98995$), so
\[
\hat{A} = \left[\frac{0.9}{0.98995},\ \frac{0.1}{0.98995},\
  \frac{0.4}{0.98995}\right] \approx [0.909,\ 0.101,\ 0.404]
\]
Check the new length:
$\sqrt{0.909^2 + 0.101^2 + 0.404^2} = \sqrt{0.826 + 0.010 + 0.163}
= \sqrt{0.999} \approx 1$ (not exactly 1 only because of rounding).
In the same way, $\hat{B} \approx [0.830,\ 0.207,\ 0.518]$ and
$\hat{C} \approx [0.110,\ 0.988,\ 0.110]$.

\textbf{Why this helps.} For vectors of length 1 the bottom of the
cosine formula is $1 \times 1 = 1$, so the cosine \emph{is} the dot
product. Check with $A$ and $B$:
\[
\hat{A} \cdot \hat{B} = 0.909 \times 0.830 + 0.101 \times 0.207 + 0.404
\times 0.518 = 0.754 + 0.021 + 0.209 = 0.984
\]
which is $\cos(A, B) \approx 0.985$ again (the last digit differs only
because the three products were rounded).

\textbf{All products at once: one matrix product.} Put the normalised
product vectors as the rows of a matrix (here 2 products, A and C; in
the exercise 500), and multiply by the normalised question vector
$\hat{B}$. Each row gives one dot product, that is one score:
\[
\begin{bmatrix} 0.909 & 0.101 & 0.404 \\ 0.110 & 0.988 & 0.110 \end{bmatrix}
\begin{bmatrix} 0.830 \\ 0.207 \\ 0.518 \end{bmatrix}
=
\begin{bmatrix}
0.909 \times 0.830 + 0.101 \times 0.207 + 0.404 \times 0.518 \\
0.110 \times 0.830 + 0.988 \times 0.207 + 0.110 \times 0.518
\end{bmatrix}
=
\begin{bmatrix} 0.984 \\ 0.353 \end{bmatrix}
\]
In NumPy this is one line; \code{@} is matrix multiplication:
\begin{lstlisting}[language=Python]
scores = index.vectors @ question   # (500, 768) @ (768,) -> (500,)
\end{lstlisting}

%--------------------------------------------------------------------
\section{Two kinds of text: queries and documents}

Questions are short (``pillow for neck pain''); product texts are long
and descriptive. The model can embed each \textbf{for its role}, set
with \code{task\_type}: the loader embedded the products as
\code{"RETRIEVAL\_DOCUMENT"}; you embed the question as
\code{"RETRIEVAL\_QUERY"}. The two kinds of vectors are made to match
each other well.

What the loader embedded for each product is stored in Firestore as
\code{embed\_text}: name, category, colour and description. For product
P1001:
\begin{note}
\small KILGAARD mattress. Category: mattress. Colour: grey. Memory foam
mattress with medium firmness, 180x200 cm. Cooling cover that keeps you
from getting too warm at night.
\end{note}
The vector only knows what is in this text. A product with a very short
description still gets a useful vector because the name and category
are added.

%--------------------------------------------------------------------
\section{768 numbers, and the length trap}\label{sec:dims}

\code{gemini-embedding-001} returns 3072 numbers by default, already of
length 1. It is trained so that the \emph{first} numbers carry the most
meaning (\textbf{Matryoshka} training, after the Russian dolls that fit
inside each other), so we can ask for only the first 768: 4 times less
storage, faster search, a small loss in quality. We use 768 because
Firestore's own vector search (Step 3) accepts at most 2048 numbers.

\textbf{The trap: shortened vectors are not length 1.} Measured on this
project:
\begin{itemize}[nosep, leftmargin=1.4em]
  \item the product name ``KILGAARD mattress'', embedded as a document,
        768 numbers: length $0.590$;
  \item the question ``a couch my guests can sleep on'', 768 numbers:
        length $0.5905$.
\end{itemize}
The loader already normalised the product vectors before storing them.
\textbf{The question vector is yours to normalise.}

\textbf{What happens if you forget.} Write $q$ for the question vector
and $\hat{q} = q / 0.5905$ for its normalised form, so $q = 0.5905 \times
\hat{q}$. For any product $\hat{p}$:
\[
\hat{p} \cdot q = \hat{p} \cdot (0.5905 \times \hat{q})
  = 0.5905 \times (\hat{p} \cdot \hat{q})
  = 0.5905 \times \cos(p, q)
\]
Every score is multiplied by the same number $0.5905$. Measured: the
best sofa bed scores $0.716$ as a cosine and $0.423$ without
normalising; check: $0.716 \times 0.5905 = 0.423$. Multiplying every
score by the same positive number does \emph{not} change their order, so
the top 5 is the same. That is why the bug hides. It breaks anything
that reads the score itself: a rule like ``only show results above
0.6'' would now show nothing, and scores of different questions can no
longer be compared.

%--------------------------------------------------------------------
\section{Filters: a mask}

The page has a category choice and a max price. A price limit is a
number to compare, not a meaning, so it is not done with vectors. It is
done with a \textbf{mask}: one True or False per product, ``may this
product be shown?''.

\textbf{Worked example.} Four products, max price 4000 kr:
\begin{center}
\begin{tabular}{@{}lrrcc@{}}
\toprule
Row & Score & Price (kr) & $\text{price} \le 4000$? & Score after the mask \\
\midrule
0 & $0.82$ & 3000 & True  & $0.82$ \\
1 & $0.91$ & 6500 & False & $-\infty$ \\
2 & $0.40$ & 1200 & True  & $0.40$ \\
3 & $0.77$ & 2500 & True  & $0.77$ \\
\bottomrule
\end{tabular}
\end{center}
Row 1 has the best score but is too expensive. Setting its score to
$-\infty$ (minus infinity, lower than any real score) pushes it to the
very end. The top 2 is now row 0 ($0.82$), then row 3 ($0.77$).

In NumPy:
\begin{lstlisting}[language=Python]
mask = index.prices <= 4000                  # [True, False, True, True]
scores = np.where(mask, scores, -np.inf)     # keep the score where True
\end{lstlisting}
Two conditions together use \code{\&} (``and''), each in brackets:
\begin{lstlisting}[language=Python]
mask = (index.categories == "sofa") & (index.prices <= 4000)
\end{lstlisting}

Two traps:
\begin{itemize}[leftmargin=1.4em]
  \item If only 2 products pass the filters and the page asks for 24,
        return 2, not 24. Otherwise the last 22 are products with score
        $-\infty$, that is products that should not be shown.
  \item Real effect, measured with the reference: for the question
        ``sofa under 4000 kr'' with no filter, all 5 results cost between
        6\,569 and 10\,269 kr. The words ``under 4000'' do not make a
        vector cheaper. With \code{max\_price=4000} the mask fixes it.
        Reading the limit out of the question automatically is Step 3.
\end{itemize}

%--------------------------------------------------------------------
\section{Putting it together: what search() does}

\code{search(query, filters, k)} in \code{api/search.py}, step by step:
\begin{enumerate}[leftmargin=1.6em, itemsep=0.25em]
  \item \code{load\_index()}: read the 500 products from Firestore once
        and keep four arrays in the same order: IDs, vectors ($500 \times
        768$), categories and prices. Row $i$ of each is product $i$.
        \code{@lru\_cache} keeps the result, so later searches do not
        read Firestore again.
  \item \code{embed\_query(query)}: 768 numbers, as
        \code{RETRIEVAL\_QUERY}.
  \item \code{normalise(...)}: the question vector to length 1
        (Section~\ref{sec:dims}).
  \item One matrix product: 500 scores (Section~3.4).
  \item \code{allowed(index, filters)}: the mask; scores of the
        others become $-\infty$ (Section~6).
  \item \code{top\_k(scores, k)}: the row numbers of the $k$ highest
        scores, best first; but never more than the number of allowed
        products.
  \item Return \code{Hit(product\_id=..., score=...)} for each row.
\end{enumerate}
The API (\code{api/app.py}) turns the hits into product cards with name,
price and photo. You do not touch it.

%--------------------------------------------------------------------
\section{How to know if the search is good}

Each of the 23 test questions in \code{data/eval/product\_queries.json}
lists the products that are correct for it. Two measures on the top 5:

\begin{itemize}[leftmargin=1.4em]
  \item \textbf{hit@5}: is at least one correct product in the top 5?
        1 or 0.
  \item \textbf{precision@5}: what share of the top 5 is correct? From 0
        to 1.
\end{itemize}

\textbf{Worked example} (real, with the reference): ``rug I can put in
the washing machine''. Of the 5 results, the first 4 are correct and the
5th is not. So hit@5 $= 1$ (at least one is correct) and precision@5 $=
4 / 5 = 0.8$. The check averages both over all questions of the same
kind.

The questions have three kinds: \code{meaning} (15), \code{language} (5,
in Danish or German) and \code{needs\_filter} (3: ``sofa under 4000
kr'', ``sofa that is not leather'', ``mattress 180x200''). The last kind
is not graded: embeddings alone cannot do ``under'', ``not'' or exact
sizes. The reference reaches hit@5 $1.00$ and precision@5 $0.99$ on
\code{meaning}, $1.00$ and $0.96$ on \code{language}. To pass, both kinds
together need hit@5 $\geq 0.95$ and precision@5 $\geq 0.90$.

%--------------------------------------------------------------------
\section{The exercise}

\begin{task}{Exercise · fill in api/search.py}
Fill in TODO 1 to TODO 6 in \code{api/search.py}, from top to bottom.
After each function, run from the repo root:
\begin{lstlisting}[backgroundcolor=\color{white}]
uv run python steps/01_embeddings/check.py
\end{lstlisting}
Part 1 of the check tests \code{normalise}, \code{allowed} and
\code{top\_k} on the numbers of this document, without calling any
service. Part 2 asks all 23 test questions.
\end{task}

\textbf{What Part 1 checks, so you can predict it by hand:}
\begin{itemize}[leftmargin=1.4em, itemsep=0.2em]
  \item \code{normalise([0.9, 0.1, 0.4])} is about \code{[0.909, 0.101,
        0.404]} (Section~3.4).
  \item \code{normalise} on the matrix with rows $[3, 4]$ and $[0, 2]$
        gives rows $[0.6, 0.8]$ and $[0, 1]$. First row: length $5$, then
        $3/5 = 0.6$ and $4/5 = 0.8$. Second row: length $2$, then $0/2 =
        0$ and $2/2 = 1$.
  \item \code{allowed} on the 4 products of Section~6 (categories sofa,
        sofa, rug, sofa): max price 4000 keeps rows 0, 2, 3; sofa and max
        price 4000 keeps rows 0 and 3; no filters keeps all four.
  \item \code{top\_k([0.2, 0.9, 0.1, 0.7, 0.5], 3)} is \code{[1, 3, 4]}:
        the scores $0.9$, $0.7$ and $0.5$ sit at positions 1, 3 and 4.
\end{itemize}

\textbf{See it on the page.} Run the API and the page locally (two
terminals):
\begin{lstlisting}[backgroundcolor=\color{white}]
uv run uvicorn api.app:app --reload --port 8000
cd web && npm run dev
\end{lstlisting}
Search ``a couch my guests can sleep on''. The label at the right shows
``your search'' instead of ``keyword placeholder''. Try the max-price
filter with ``sofa under 4000 kr''.

\begin{note}
\textbf{Live site.} The deployed API runs as the service account
\code{nordhjem-api}, which has no permissions yet. Before your search
can run on \code{nordhjemstore.web.app}, that account needs to read
Firestore and call Vertex AI (roles \code{datastore.user} and
\code{aiplatform.user}). We do that together when you deploy.
\end{note}

%--------------------------------------------------------------------
\section{Experiments, after it passes}

Change one thing at a time, run the check again, and write the numbers
in a comment at the top of \code{api/search.py}.

\begin{enumerate}[leftmargin=1.6em, itemsep=0.35em]
  \item \textbf{Forget to normalise} the question on purpose. Do hit@5
        and precision@5 change? Do the scores? Explain with
        Section~\ref{sec:dims}.
  \item \textbf{Task type.} Embed the question as
        \code{RETRIEVAL\_DOCUMENT}. What happens to the scores, and to
        the ranking?
  \item \textbf{Colour variants.} ``a couch my guests can sleep on''
        returns the same HALDAGAARD sofa in grey, beige and dark green.
        How would you show 5 \emph{different} products? (Hint: group by
        name, keep the best score.)
  \item \textbf{The limits.} Look at what ``sofa that is not leather''
        returns. Why can a vector not understand ``not''? Write your
        explanation; we discuss it before Step 3.
\end{enumerate}

\begin{note}
The reference solution is \code{steps/01\_embeddings/solution.py}. Open
it only after trying, or when stuck on one function.
\code{check.py --solution} runs the checks on it.
\end{note}

\end{document}
